\SourceLeaf{370}%
Ein solches aufrichtiges und unumwundenes Eingeständnis scheint uns der Würde der Wissenschaft angemessener als das eitle Bemühen, die Lücke, die man nicht schließen kann, hinter einem unhaltbaren Geflecht von Scheinbeweisen zu verbergen.

Der Verfasser der ersten Schrift hatte bereits vor 15 Jahren in einer kleinen Abhandlung, \textit{Tentamen novae parallelarum theoriae notione situs fundatae}, einen ähnlichen Versuch unternommen, indem er alles auf den Begriff der Identität der Lage zu stützen suchte. Er definiert Parallelen als Geraden, die dieselbe Lage haben. Daraus schließt er, dass jede dritte Gerade sie notwendig unter gleichen Winkeln schneiden muss, weil diese Winkel lediglich das Maß dafür seien, wie sehr sich die Lage dieser dritten Geraden von den Lagen der beiden Parallelen unterscheidet. Diese Beweisart wird in der vorliegenden neuen Schrift wiederholt, und wir könnten kaum sagen, dass sie durch die eingeflochtenen philosophischen Betrachtungen an Stärke gewonnen hätte. Der Behauptung auf S.~24, Notionem situs e geometria adeo non excludi posse, ut potius notionibus eius fundamentalibus annumeranda sit, \textit{dudum omnes agnovere geometrae}, muss jeder Geometer in dem Sinn widersprechen, in dem der Verfasser den Begriff der Lage in seinem Beweis gebraucht. Gehen wir von der Definition des Verfassers aus, Situs est modus, quo plura coëxistunt vel iuxta se existunt in spatio, so ist Lage ein reiner Verhältnisbegriff. Man kann durchaus sagen, dass zwei Geraden \(A,B\) eine bestimmte Lage zueinander haben, die mit der gegenseitigen Lage zweier anderer Geraden \(C,D\) übereinstimmt. In seinem Beweis verwendet der Verfasser das Wort Lage jedoch als absoluten Begriff, indem er von der Identität der Lage zweier Geraden spricht, die nicht zusammenfallen. Diese Bedeutung ist offenkundig leer und haltlos, solange wir nicht wissen, was wir unter einer solchen Identität verstehen und woran wir sie erkennen sollen. Soll man sie an der Gleichheit der Winkel mit \textit{einer} dritten Geraden erkennen, dann wissen wir ohne vorherigen Beweis noch nicht, ob dieselbe Gleichheit auch für die Winkel mit einer vierten Geraden gilt. Soll die Gleichheit der Winkel mit \textit{jeder} anderen Geraden das Kriterium sein, dann wissen wir wiederum nicht, ob gleiche Lage ohne Zusammenfallen möglich ist. Wir stehen also \textit{nach} dem Beweis des Verfassers genau da, wo wir \textit{vor} ihm standen.

Ein großer Teil der Schrift befasst sich mit der gegen \textsc{Kant} gerichteten Behauptung,
